% Lösungen Einheit 3 von 5 – Prisma (zusammenbau v0.1)
\einheitenkopf{Einheit 3 von 5 · Prisma}

\erg{20}{a) Grundfläche: eines der beiden gleichen Dreiecke (die Deckfläche ist das andere); Höhe: eine Kante zwischen den Dreiecken, der Abstand von Grund- und Deckfläche – auch wenn das Prisma liegt \quad b) $18 \cdot 7 = 126$ cm³ \quad c) Grundfläche $7 \cdot 3 = 21$ cm²; $21 \cdot 5 = 105$ cm³ \quad d) Grundfläche $8 \cdot 5 : 2 = 20$ cm²; $20 \cdot 9 = 180$ cm³ \quad e) Grundfläche $(7 + 3) : 2 \cdot 4 = 20$ cm²; $20 \cdot 6 = 120$ cm³ \quad f) Grundfläche $2{,}4 \cdot 1{,}6 : 2 = 1{,}92$ m²; $1{,}92 \cdot 3 = 5{,}76$ m³ \quad g) Umfang $4 + 7 + 8 = 19$ cm; $19 \cdot 5 = 95$ cm² \quad h) Grundfläche $3 \cdot 4 : 2 = 6$ cm²; Mantel $(3 + 4 + 5) \cdot 7 = 84$ cm²; $2 \cdot 6 + 84 = 96$ cm² \quad i) Es fehlen die zwei Dreiecke mit den Seiten $4$ cm, $5$ cm und $6$ cm, je eins oben und unten an einem Rechteck. \quad j) $210 : 35 = 6$ cm \quad k) Deckfläche $\tfrac{1}{2} \cdot 1{,}8 \cdot 1{,}56 = 1{,}404$ m²; Volumen $1{,}404 \cdot 0{,}75 \approx 1{,}05 \approx 1{,}1$ m³}
\erg{21}{a) Beim Dreieck fehlt das Teilen durch zwei. Richtig: Grundfläche $7 \cdot 6 : 2 = 21$ cm², Volumen $21 \cdot 5 = 105$ cm³ \quad b) Das Prisma ist ein Stapel gleicher Schichten in der Form der Grundfläche. Eine Schicht von einem Zentimeter Dicke hat so viele cm³, wie die Grundfläche cm² hat; es gibt so viele Schichten, wie das Prisma in cm hoch ist. \quad c) $9 \cdot 3{,}5 : 2 = 15{,}75$ m²; $15{,}75 \cdot 11 = 173{,}25$ m³ – ja, weniger als $180$ m³ \quad d) zwei Dreiecke mit den Seiten $8$ cm, $5$ cm, $5$ cm; drei Rechtecke $8$ cm × $7$ cm, $5$ cm × $7$ cm, $5$ cm × $7$ cm nebeneinander, die Dreiecke oben und unten am Rechteck}
